% ---------------------------------------------------------------------
\section{Results}\label{sec:results}
% ---------------------------------------------------------------------
% SKELETON. No experiment has been run yet. Every quantitative statement
% is a \RES{} placeholder that asks only for values produced by
% analysis/scripts/make_all.py (tables T1-T7, T3b, TS1-TS3; figures fig3-fig6,
% figS1-figS6; stats/tests_T*.csv). Generated tables and figures appear
% automatically once they exist under \analysisdir (default ../../analysis)
% or as local copies in tables/ and figures/.

\TODO{Skeleton: the authors replace every sentence frame below with their own text, fill the red placeholders from the analysis outputs and write all interpretation themselves}.

\subsection{Main comparison}\label{sec:res-main}

\Tab{tab:t1} compares all methods at the training sizes $N\in\{10,20\}$, pooled over profiles and ratings.
At these sizes, the gap of \method{} is \RES{gap and IQM with 95\% confidence interval, \Tab{tab:t1}}.
Relative to the other learned methods, the difference in penalised profit is \RES{differences, confidence intervals and Holm-adjusted $p$-values, stats/tests\_T1.csv}.
Among the rules and model-predictive methods, \RES{best rule with the safety layer and the gaps of eMPC-HiGHS, \Tab{tab:t1}}.

\restable{T1}{%
\begin{table}[t]
\centering
\caption{Main comparison at the training station sizes $N\in\{10,20\}$ (three arrival profiles and both tightness levels pooled; placeholder until the analysis scripts have run).}
\label{tab:t1}
\footnotesize
\setlength{\tabcolsep}{3.5pt}
\begin{tabular}{@{}lcccccc@{}}
\toprule
Method & $J^{\sm{pen}}$ [\texteuro] & Gap [\%] & IQM gap [\%] & Profit [\texteuro] & Not served [kWh] & Overload [kWh] \\
\midrule
\method{} (ours) & \RES{} & \RES{} & \RES{} & \RES{} & \RES{} & \RES{} \\
Attention+penalty & \RES{} & \RES{} & \RES{} & \RES{} & \RES{} & \RES{} \\
Attention+PPO-Lagrangian & \RES{} & \RES{} & \RES{} & \RES{} & \RES{} & \RES{} \\
Deep Sets+safety layer & \RES{} & \RES{} & \RES{} & \RES{} & \RES{} & \RES{} \\
Rules and eMPC-HiGHS & \RES{} & \RES{} & \RES{} & \RES{} & \RES{} & \RES{} \\
\bottomrule
\end{tabular}
\end{table}}

\subsection{Zero-shot transfer across station sizes}\label{sec:res-scaling}

\Fig{fig:zeroshot} plots the penalised profit per charger and the gap against $N$ for weights trained at $N\in\{10,20\}$.
At $N=50$, $100$ and $200$, the gap of \method{} is \RES{gap per $N$ with 95\% confidence interval, \Tab{S-tab:t2}}, and its transfer gap to the reference trained at $N=50$ is \RES{\Tab{S-tab:t2}}.
EV-GNN reaches \RES{gap in its training configuration and zero-shot, \Tab{S-tab:t2}}, and at the stress level $\rho_{\sm{tr}}=3$ kW the gap of \method{} is \RES{\Tab{S-tab:ts2}}.

\begin{figure}[t]
\centering
\resultfig{fig3_zero_shot_scaling.pdf}{Penalised profit per charger and optimality gap versus $N$ per method, 95\% confidence intervals; training sizes marked.}
\caption{Zero-shot transfer across station sizes for $N$ from 5 to 200.}
\label{fig:zeroshot}
\end{figure}

\subsection{Constraint satisfaction}\label{sec:res-constraints}

\Fig{fig:constraints} shows energy not served, overload energy and the share of EVs without infeasible steps (E4) against $N$ for the main methods.
With the safety layer, the overload energy is \RES{value from \Tab{S-tab:t3}; Theorem~\ref{thm:invariance}(a) predicts zero}, and without it the methods leave \RES{energy not served and overload, \Tab{S-tab:t3}}.
E4 of \method{} is \RES{per $N$, \Tab{S-tab:t3b}}, and the aggregate limits bind on \RES{binding rates per $N$, \Tab{S-tab:t3b}}.
Infeasible steps fall into the fallback cases as \RES{counts per case, \Tab{S-tab:t3b}}, and the shortfall reaches \RES{ratio to the deficit bound, \Tab{S-tab:t3}}.

\begin{figure}[t]
\centering
\resultfig{fig4_constraint_satisfaction.pdf}{Energy not served, overload energy and E4 versus $N$ for the main methods.}
\caption{Constraint satisfaction across station sizes, including the share of EVs whose stay contains no infeasible step (E4).}
\label{fig:constraints}
\end{figure}

\subsection{Cost of the guarantee}\label{sec:res-cost}

\Fig{fig:cost} plots the cost of the guarantee of \Sec{sec:metrics} against $N$.
In penalised profit, Attention+penalty minus \method{} is \RES{$\Delta J^{\sm{pen}}$ in \texteuro{} per EV served and in \% of $|J^{\sm{LP,ceil}}|$ per $N$, \Tab{S-tab:t5}}.
Measured in raw profit, the difference is \RES{$\Delta$ profit per $N$, \Tab{S-tab:t5}}, with the tests in \RES{stats/tests\_T5.csv}.

\begin{figure}[t]
\centering
\resultfig{fig6_cost_of_guarantee.pdf}{Cost of the guarantee ($\Delta J^{\sm{pen}}$ and $\Delta$ profit, \texteuro{} per EV served) versus $N$.}
\caption{Cost of the guarantee across station sizes. A positive value means that Attention+penalty scored higher than \method{} on the same episodes.}
\label{fig:cost}
\end{figure}

\subsection{Further results}\label{sec:res-further}

In \Tab{S-tab:t4}, the variants without departure bound, with the layer at deployment only and without admission control give \RES{gaps relative to \method{}, \Tab{S-tab:t4}}.
For the conformal variant, the gap and the episode-averaged coverage are \RES{per $\alpha$, \Tab{S-tab:t7}}, and one control step at $N=200$ takes \RES{inference time, \Tab{S-tab:t6}}.
The sensitivity studies give \RES{\Tab{S-tab:ts2}}, and Figure~\ref{S-fig:exampleday} shows an example episode.
